\clearpage
\phantomsection% 
\addcontentsline{toc}{chapter}{ABSTRAK}
\thispagestyle{fancy}

\begin{center}
	\large \bfseries \MakeUppercase{Abstrak}\\
	\normalsize \normalfont {\thetitle}\\
	\normalsize \normalfont {\theauthor}\\
	\bigskip
	
	\normalsize \normalfont \justifying \singlespacing
	Pembelajaran kolaboratif merupakan komponen vital dalam pendidikan modern untuk mengembangkan keterampilan interpersonal dan pemecahan masalah mahasiswa. Namun, pembentukan kelompok manual sering menghasilkan ketidakseimbangan keterampilan, konflik interpersonal, dan partisipasi yang tidak merata. Penelitian ini mengembangkan sistem informasi EduTeams berbasis kecerdasan buatan untuk otomatisasi pembentukan kelompok mahasiswa menggunakan API Edu2Com yang mempertimbangkan kepribadian MBTI, keterampilan, preferensi topik, dan keseimbangan gender. Pengembangan sistem menggunakan metodologi Agile Kanban untuk memastikan fleksibilitas dan kontinuitas perbaikan, sementara pengujian menggunakan metode grey-box testing. Hasil penelitian menunjukkan bahwa mayoritas fitur telah berhasil diimplementasikan dan diuji. Pengujian komprehensif mencapai tingkat keberhasilan yang tinggi dengan cakupan kode yang memadai. Integrasi Edu2Com API berhasil menghasilkan pembentukan kelompok otomatis sesuai ekspektasi. Pengujian non-fungsional pada tiga browser (Chromium, Firefox, Mobile Chrome) menunjukkan portabilitas dan dukungan multi-bahasa yang baik, serta skor Lighthouse yang memenuhi standar. Kegagalan minor terjadi pada interaksi tertentu di perangkat mobile. Sistem EduTeams telah mencapai tingkat kematangan yang baik dan siap untuk penggunaan dalam lingkungan produksi dengan dukungan multi-bahasa (Indonesia dan Inggris), keamanan berbasis role-based access control, dan performa yang memenuhi standar kualitas perangkat lunak modern.\par

    \vspace{20pt}
	\raggedright \textbf{Kata Kunci: pembelajaran kolaboratif, pembentukan kelompok otomatis, kecerdasan buatan, Agile Kanban, grey-box testing}
	
	\vfill
	
\end{center}
\clearpage